\chapter{Continuum limit and Symanzik's Improvement program}

The Taylor expansion of the simple Wilson action for the gauge and fermion
degrees of freedom showed the presence of higher dimensional operators.
Since these are suppressed by powers of the lattice spacing, they are
irrelevant operators in the language of the renormalization group, i.e.\ they
vanish at the fixed point at $a \to 0$. Another way of expressing the same
thing is as follows. The lattice theory with a hard cut-off at $\pi/a$ can be
regarded as an effective theory [49]. Integration of momenta from
$\infty \to \pi/a$ generates effective interactions. Thus, in general, we can
write the action as
\begin{equation}
S_C = S_L(\pi/a) + \sum_n \sum_i a^n C_i^n A_i^n(\pi/a) \,,
\label{eq:8.1}
\end{equation}
where $S_C$ ($S_L$) is the continuum (lattice) action, and the first sum is
over all operators of a given dimension, and the second over all dimensions
greater than four. Similarly, the operators used to probe the physics can be
written as
\begin{equation}
O_C = O_L(\pi/a) + \sum_n \sum_i a^n D_i^n O_i^n(\pi/a) \,,
\label{eq:8.2}
\end{equation}
where in this case quantum corrections can induce mixing with lower dimension
operators also. The coefficients $C_i^n$ and $D_i^n$ are both functions of the
coupling $\alpha_s$. Parenthetically, one important reason for preserving
gauge, Lorentz, and chiral symmetries at finite $a$ is that they greatly
restrict the set of possible operators of a given dimension.

Now consider the expectation value of a given operator
\begin{equation}
\begin{split}
\langle O_C \rangle &= \int dU\, O_C e^{-S_C} \\
&= \int dU \left( O_L + \sum_n \sum_i a^n D_i^n O_i^n \right)
   e^{-S_L - \sum_m \sum_i a^m C_i^m A_i^m} \\
&= \int dU \left( O_L + \sum_n \sum_i a^n D_i^n O_i^n \right)
   \left( 1 - \sum_m \sum_i a^m C_i^m A_i^m + \cdots \right) e^{-S_L} \\
&= \int dU\, O_L e^{-S_L}
   + \int dU\, O_L \left( \sum_m \sum_i a^m C_i^m A_i^m + \cdots \right) e^{-S_L} \\
&\qquad + \int dU \left( \sum_n \sum_i a^n D_i^n O_i^n \right)
   \left( 1 - \sum_m \sum_i a^m C_i^m A_i^m + \cdots \right) e^{-S_L} \,.
\end{split}
\label{eq:8.3}
\end{equation}
Note that the operators $A_i^n$, which appear in the action, are summed over
all space-time points. This leads to contact terms in Eq.~\ref{eq:8.3} which
I shall, for brevity, ignore as they necessitate a proper definition of the
operators but do not change the conclusions.

The basis for believing that LQCD provides reliable results from simulations
at finite $a$ is that the contributions of the lattice artifacts (last two
terms in the last expression) are small, can be estimated, and thus removed.
At this point I will only make a few general statements about the calculation
of expectation values which will be embellished on in later lectures.
\begin{itemize}
\item The contribution of the last two terms in Eq.~\ref{eq:8.3} vanishes in
the limit $a \to 0$ unless there is mixing with operators of lower dimension.
In such cases the mixing coefficient has to be determined very accurately
(necessitating non-perturbative methods) otherwise the $a \to 0$ limit is
divergent.
\item Calculations have to be done at values of $a$ at which contributions of
irrelevant terms organized in successive powers in $a$ decrease at least
geometrically.
\item To achieve continuum results simulations are carried out at a sufficient
number of values of $a$ to allow reliable extrapolation to $a = 0$. The
``fineness'' of the lattice spacing $a$ at which to do the calculations and
the number of points needed depends on the size of the corrections and how
well we can determine the functional form to use in the extrapolation.
\item To improve the result to any given order in $a$, one has to improve both
the action and the operators to the same order. Such an approach, based on
removing lattice artifacts organized as a power series expansion in $a$, is
called the Symanzik improvement program.
\end{itemize}

Eq.~\ref{eq:8.3} shows clearly that there is fair degree of flexibility in
constructing LQCD. We are free to add any irrelevant operator with a sensible
strength and still recover QCD in the continuum limit. The only difference
between various constructions of the action is the approach to the continuum
limit. What we would like to do is improve the action and operators so as to
get continuum results on as coarse lattices as possible. On the other hand
improvement requires adding irrelevant terms to both the action and the
operators. This increases the complexity of the calculations and the
simulation time. Thus, the strategy for minimizing discretization errors is a
compromise -- optimize between simplicity of the action and operators
(defined in terms of the cost of simulation, code implementation, and
analyses) and the reduction of the discretization errors evaluated at some
given fixed value of the lattice spacing, say $a \approx 0.1$ fermi. The
remaining errors, presumably small, can then be removed by extrapolating the
results obtained from simulations at a few values of $a$ to $a = 0$.

The number of floating points operations needed in the simulations of LQCD
scale roughly as $L^6$ in the quenched approximation and $L^{10}$ with
dynamical fermions. Current lattice simulations of QCD are done for lattice
spacing $a$ varying in the range $2\,\mathrm{GeV} \leq a^{-1} \leq
5\,\mathrm{GeV}$. For these parameter values the corrections due to the
$O(a)$, $O(ma)$, $O(pa)$, $O(a^2)$, \ldots\ terms in the Wilson action are
found to be large in many observables. The importance of improving the
lattice formulation (action and operators) is therefore self-evident,
especially for full QCD. Thus, the search for such actions is a hot topic
right now. Before discussing these let me first summarize what one hopes to
gain by, and the criteria by which to judge, such an improvement program.
\begin{enumerate}
\item Improved scaling, i.e.\ discretization errors are reduced. Consequently
one can work on coarser lattices (smaller $L = N_S$) for a given accuracy.
Even a factor of 2 reduction in $L$ translates into $2^{10} \approx 1000$ in
computer time!

\item Better restoration of rotational and internal symmetries like chiral,
staggered flavor, etc..

\item The trajectory along which the continuum limit is taken should not pass
close to extraneous phase transitions that are artifacts of the lattice
discretization (see Section 15 for a discussion of phase transitions in the
lattice theory). The problem with such transitions is that in their vicinity
the desired continuum scaling behavior can be modified for certain
observables, i.e.\ the artifact terms in Eq.~\ref{eq:8.3} may not have a
geometric convergence in $a$ and may thus be large.

\item The study of topology on the lattice is particularly sensitive to short
distance fluctuations that are artifacts of discretizing the gauge action.
For the Wilson gauge action one finds that the topological charge, calculated
using L\"uscher's geometric method [50], jumps from $0 \to 1$ at a core
radius of the instanton of $\rho_c a \approx 0.7$, whereas the action of such
instantons, $S_{\rho_c}$, is significantly less than the classical value
$S^{\mathrm{cl}} = 8\pi^2/\beta$. The entropy factor for such small
instantons with $S < S^{\mathrm{cl}}$ overwhelms the Boltzmann suppression,
leading to a divergent topological susceptibility in the continuum limit.
These short distance artifacts are called dislocations and it is not
a~posteriori obvious how to separate the physical from the unphysical when
$\rho \approx \rho_c$. It is, therefore, very important to start with a
lattice formulation in which such artifacts are excluded, and there exist
only physical instantons, i.e.\ those with mean core size that has the
correct scaling behavior with $g$. One finds that dislocations are suppressed
by improved actions, i.e.\ improving the gauge action increases $S_{\rho_c}$
[51] and better definition of the topological charge [52].

\item Lattice perturbation theory becomes increasingly complicated as more
and more irrelevant operators are added to the action. Even though, with the
advent of non-perturbative methods for determining renormalization constants,
the reliance on lattice perturbation theory has decreased, perturbative
analyses serve as a valuable guide and check, so we would like to retain this
possibility. What one expects for an improved action is that the lattice and
continuum value of $\alpha_s$, and the coefficients of the series expansion
in $\alpha_s$ for a given observable, are much closer. In that case a much
more reliable matching between the two theories is possible. At present the
1-loop matching between the lattice and continuum theories is one of the
larger sources of uncertainties.

\item The Adler-Bell-Jackiw (ABJ) Anomaly: In the continuum the flavor
singlet axial current has an anomaly. This shows up as part of the UV
regularization scheme. In the Pauli-Villars scheme one introduces a heavy
mass as the cutoff. This breaks $\gamma_5$ invariance, while in dimensional
regularization, $\gamma_5$ is not well-defined away from 4 dimensions.
Naive fermions, as shown below, preserve the $\gamma_5$ invariance but at the
expense of flavor doubling. A consequence of the doublers is a pairwise
cancellation and no net anomaly. A viable lattice theory should reproduce the
ABJ anomaly.
\end{enumerate}

Before discussing such ``improved'' actions, it is instructive to revisit the
construction of the naive Dirac action to highlight its problems.
